\section*{Заключение}
\addcontentsline{toc}{section}{Заключение}

Проведенный комплексный теоретико-методологический и исторический анализ позволяет сформулировать следующие ключевые выводы:

\begin{enumerate}
    \item Политическая культура представляет собой многоуровневую систему ментальных, эмоциональных, ценностных и поведенческих ориентаций, которая выступает важнейшим регулятором и стабилизатором макрополитической системы. Ее когнитивные, аффективные и оценочные компоненты определяют границы легитимности институтов публичной власти.
    \item Классическая концепция Г. Алмонда и С. Вербы доказала, что устойчивость демократического транзита гарантируется не чистым активизмом, а формированием смешанной «гражданской культуры». В ней рациональное участие сбалансировано уважением к закону и традиционным конформизмом, что удерживает систему в равновесии.
    \item В контексте теории С. Хантингтона эволюция политических ценностей носит нелинейный характер. Опережающий рост участия масс при слабой укорененности культуры компромисса ведет к политическому упадку, доказывая, что культурные институты обладают колоссальной инерцией и способны реанимировать авторитарные паттерны в эпохи кризисов.
    \item Российская политическая культура характеризуется глубокими, устойчивыми историческими доминантами, такими как этатизм, патернализм, персонификация власти и ценностный дуализм. На современном этапе отечественная ментальность демонстрирует сложный синтез традиционных государственнических архетипов с постепенно нарастающими элементами рационально-субъектного, цифрового участия, что определяет контуры долгосрочной эволюции всей политической системы страны.
\end{enumerate}
